problem:  
Let \( (M^n,g) \) be a closed Riemannian manifold with sectional curvature \( K_g \ge 1 \).  Assume that  
\[
\frac{\operatorname{Vol}(M,g)}{\operatorname{Vol}(S^n)} = 1 - \varepsilon,
\]
for small \(\varepsilon>0\).  Grove–Petersen’s theorem ensures spherical *homeomorphism* when the volume is sufficiently close to maximal.  

**Question (Quantitative smooth rigidity at almost maximal volume):**  
Determine whether there exist constants \( \varepsilon_n>0 \) and \( \alpha_n>0\) such that for all \( M \) with \( K_g \ge 1 \) and \( \varepsilon < \varepsilon_n \), one can find a diffeomorphism \( \Phi:M\to S^n \) with  
\[
\|\Phi^* g_S - g\|_{C^{2}} \le C_n\, \varepsilon^{\alpha_n}.
\]
In other words, can we quantify *how smoothly* a positively curved manifold with almost maximal volume approximates the round sphere?  Is there a universal Hölder or Lipschitz rate \( \alpha_n \) of “metric rigidity” near maximal volume?  

Equivalently, does near-maximal volume imply **quantitative \(C^k\)-closeness** (for some fixed \(k\)) to the round metric up to diffeomorphism?  

Why is it a "good" problem:  
This question pushes beyond Grove–Petersen’s *topological* almost-rigidity by asking for a **quantitative, smooth geometric stability** statement.  It probes whether “almost maximal volume” under a lower sectional curvature bound controls not only topology, but also explicit metric closeness to the canonical model.  

It is “good” because:  

1. **It strengthens a classical boundary.**  
   Rather than re‑stating known topological results, it seeks a **rate of convergence** of the metric itself—a refined, analytic version of volume rigidity.  

2. **It bridges geometry and analysis.**  
   Establishing \(C^k\)-closeness requires coupling comparison geometry with analytic deformation theory (e.g., Ricci flow stability, eigenvalue pinching, harmonic coordinate estimates).  This would link coarse volume comparison to fine regularity analysis.  

3. **It touches the smooth vs. topological interface.**  
   A positive result implies that almost maximal volume precludes exotic smooth structures, offering a new metric pathway to understand smooth classification of positively curved manifolds.  

4. **It exhibits clarity and depth.**  
   The statement uses only curvature, volume, and smooth norm closeness—simple geometric data—but its resolution would require understanding how local metric geometry encodes global smooth topology.  

In summary, the problem transforms known qualitative results into a *quantitative geometric stability program*, connecting comparison geometry, geometric analysis, and smooth topology—a concise yet profoundly deep open question.